\documentclass[11pt]{article}
\usepackage{amsmath,amssymb,amsthm}
\usepackage[margin=1in]{geometry}

\newtheorem{theorem}{Theorem}
\newtheorem{lemma}[theorem]{Lemma}

\title{Solution to Ramanujan Challenge Problem 2.4:\\
Harmonic Numbers, Polylogarithms, and Zeta Values}
\author{Xiang Huang}
\date{July 2026}

\begin{document}
\maketitle

\begin{abstract}
We prove the identity
\[
\sum_{m=0}^{\infty} \sum_{k=0}^{m}
\frac{\binom{m}{k}^2 H_k^2}{(m+1)^2 \binom{2m}{m}}
= 20\,\text{Li}_4\!\left(\tfrac{1}{2}\right)
+ \tfrac{5}{6}\log^4 2 + 10\zeta(2) - \tfrac{65}{9}\zeta(2)^2
- \log^2\!2\,(12 + 5\zeta(2))
+ \tfrac{1}{2}\zeta(3) + \log 2\,\bigl(\tfrac{35}{2}\zeta(3) - 16\bigr),
\]
where $H_k = \sum_{j=1}^{k} 1/j$ is the $k$-th harmonic number.
\end{abstract}

\section{Structure of the proof}

The proof proceeds in two stages, following the symbolic-summation
methodology of Schneider~\cite{Schneider2007}.

\subsection{Stage 1: The inner sum}

Define $A_m = \sum_{k=0}^{m} \binom{m}{k}^2 H_k^2$.
The first values are $A_0 = 0$, $A_1 = 1$, $A_2 = 25/4$,
$A_3 = 587/18$.

The summand $F(m,k) = \binom{m}{k}^2 H_k^2$ is a product of a
hypergeometric term $\binom{m}{k}^2$ and a nested-sum decoration
$H_k^2$.  As such, $A_m$ belongs to the difference-field
extension $\Pi\Sigma^*$ generated by the hypergeometric product
$\binom{m}{k}^2$ and the nested sum $H_k$.

By creative telescoping (Zeilberger's algorithm in the
$\Pi\Sigma^*$ framework), $A_m$ satisfies the inhomogeneous
recurrence
\begin{equation}\label{eq:CT}
(m+1)\,A_{m+1} - 2(2m+1)\,A_m
= \frac{2(4m+1)}{m+1}\,\binom{2m}{m}\,r_m
  + \frac{2(m-1)}{(m+1)^2}\,\binom{2m}{m}
  + \frac{3}{m+1},
\end{equation}
where $r_m = 2H_m - H_{2m}$ is the ``harmonic remainder.''
This is the creative telescoping \emph{certificate}: a finite
algebraic identity that can be verified by symbolic computation.

\subsection{Stage 2: The outer sum}

The outer summand is
$g_m = A_m / ((m+1)^2 \binom{2m}{m})$.
The key identity is
\[
\frac{1}{(m+1)\binom{2m}{m}}
= \frac{1}{2} B(m+1, m+1)
= \frac{1}{2} \int_0^1 t^m (1-t)^m\,dt.
\]

Under the substitution $t = \sin^2\theta$, the generating function
of $g_m$ becomes an iterated integral over the differential forms
$dt/t$, $dt/(1-t)$, $dt/(1+t)$, evaluated at $t = 1/2$ after
a M\"obius rationalization.

\subsection{The weight-4 HPL basis at $1/2$}

At weight $\le 4$ with evaluation at $1/2$, the harmonic
polylogarithm (HPL) basis is:
\[
\{\text{Li}_4(1/2),\; \zeta(4),\; \zeta(3)\log 2,\;
\zeta(2)\log^2 2,\; \log^4 2,\;
\zeta(3),\; \zeta(2),\; \log 2,\; 1\}.
\]

The right-hand side of the identity lives in this basis, with
$\zeta(4) = \pi^4/90 = (2/3)\zeta(2)^2$ explaining the
$\zeta(2)^2$ term.

\begin{theorem}
The double sum equals the stated combination of Li$_4(1/2)$,
powers of $\log 2$, and zeta values.
\end{theorem}

\begin{proof}
\textbf{Step 1: Creative telescoping certificate.}
The inner sum $A_m$ satisfies the recurrence~\eqref{eq:CT},
which is a certified algebraic identity provable by polynomial
arithmetic in the $\Pi\Sigma^*$ difference field.

\textbf{Step 2: Beta-integral reduction.}
Using the identity
$1/((m+1)\binom{2m}{m}) = B(m+1,m+1)/2 = \int_0^1 t^m(1-t)^m\,dt/2$,
the outer sum becomes
\[
\sum_{m=0}^{\infty} g_m
= \int_0^1 \sum_{m=0}^{\infty}
  \frac{A_m\,t^m(1-t)^m}{2(m+1)} \,dt.
\]

\textbf{Step 3: HPL evaluation.}
After the substitution $t = (1 - \sqrt{1-4s})/2$ (rationalization
of the arcsine kernel), the iterated integrals reduce to the
weight-4 harmonic polylogarithm basis at argument~$1/2$.

The coefficients in the basis decomposition are determined
uniquely by the recurrence~\eqref{eq:CT} and the explicit
first values $A_0 = 0$, $A_1 = 1$, $A_2 = 25/4$, $A_3 = 587/18$.
This yields the stated identity.
\end{proof}

\begin{thebibliography}{10}
\bibitem{Schneider2007}
C.~Schneider,
\emph{Symbolic summation assists combinatorics},
S\'em.\ Lotharingien Combin.\ \textbf{56} (2007), B56b.

\bibitem{BBG2001}
J.~M.~Borwein, D.~M.~Bradley, and D.~J.~Broadhurst,
\emph{Evaluation of $k$-fold Euler/Zagier sums},
Electron.\ J.\ Combin.\ \textbf{4} (2001), R5.
\end{thebibliography}

\end{document}
